\documentclass[12pt, a4paper]{article}

% ---- Ajustes de Idioma e Fontes Modernas -------------------
\usepackage[utf8]{inputenc}
\usepackage[T1]{fontenc}
\usepackage[brazil]{babel}
\usepackage{mathpazo}
\usepackage[scaled=0.95]{helvet}
\usepackage{courier}

% ---- Pacotes de Layout e Design ----------------------------
\usepackage{geometry}
\geometry{a4paper, top=3cm, bottom=2cm, left=3cm, right=2cm}
\usepackage{setspace}
\onehalfspacing
\usepackage{indentfirst}
\usepackage{microtype}

% ---- Cores e Hiperlinks ------------------------------------
\usepackage{xcolor}
\definecolor{brandblue}{RGB}{44, 62, 80}
\definecolor{softgray}{RGB}{245, 245, 245}
\usepackage[colorlinks=true, allcolors=brandblue]{hyperref}

% ---- Elementos Matemáticos e Gráficos ----------------------
\usepackage{amsmath, amssymb, amsfonts}
\usepackage{graphicx}
\usepackage{booktabs}
\usepackage{caption}
\usepackage{subcaption}
\usepackage{float}

% ---- Customização de Títulos -------------------------------
\usepackage{titlesec}
\titleformat{\section}{\color{brandblue}\normalfont\Large\bfseries}{\thesection}{1em}{}
\titleformat{\subsection}{\color{brandblue}\normalfont\large\bfseries}{\thesubsection}{1em}{}

% =============================================================
\begin{document}
% =============================================================

% ---- CAPA ---------------------------------------------------
\begin{titlepage}
    \begin{flushleft}
        \Large \textbf{Instituto Federal do Ceará} \\
        \large Mestrado em Ciência da Computação \\
        \small Disciplina: Reconhecimento de Padrões
    \end{flushleft}
    \vspace{4cm}
    \begin{center}
        \hrule\vspace{0.5cm}
        {\huge \textbf{\color{brandblue} Classificadores Bayesianos}}\\
        \vspace{0.2cm}
        {\Large Trabalho 04: Discriminantes Linear e Quadrático}
        \vspace{0.5cm}\hrule
    \end{center}
    \vspace{3cm}
    \begin{flushright}
        \large \textbf{Raquel Maciel Coelho de Sousa}\\
        \small Professor: Dr. Ajalmar Rocha Neto
    \end{flushright}
    \vfill
    \begin{center}\small Fortaleza, \today \end{center}
\end{titlepage}

\newpage
{\hypersetup{linkcolor=black}\tableofcontents}
\newpage

% =============================================================
\section{Introdução}
% =============================================================

Este relatório descreve a implementação e avaliação do
\textbf{Classificador Bayesiano Gaussiano} com base nos
\textbf{discriminantes linear (LDA)} e \textbf{quadrático (QDA)},
desenvolvido como Trabalho~04 da disciplina de Reconhecimento de
Padrões.

Diferentemente dos trabalhos anteriores, aqui a fronteira de decisão
é derivada diretamente das funções discriminantes log-posteriori,
permitindo verificar quatro situações distintas para a matriz de
covariância: QDA com matriz completa por classe, QDA com matriz
diagonal por classe, LDA com matriz \textit{pooled} compartilhada e
LDA com matriz esférica $\Sigma = \sigma^2 I$.

O modelo foi treinado e testado sobre os cinco conjuntos de dados:
Flor~Íris, Coluna Vertebral, Breast Cancer Wisconsin, Dermatology e
Artificial~I, e seus resultados foram comparados com os
classificadores previamente implementados (Bayesiano Gaussiano
Multivariado, Naive Bayes, KNN e DMC).

% =============================================================
\section{Fundamentação Teórica}
% =============================================================

\subsection{Classificador Bayesiano Gaussiano}

O classificador Bayesiano ótimo atribui uma amostra $\mathbf{x}$ à
classe $W_i$ que maximiza a posteriori
$P(W_i \mid \mathbf{x}) \propto p(\mathbf{x} \mid W_i)\,P(W_i)$.
Quando cada verossimilhança $p(\mathbf{x} \mid W_i)$ é uma gaussiana
multivariada, é preferível trabalhar com o logaritmo natural:

\begin{equation}
    g_i(\mathbf{x}) = \ln p(\mathbf{x} \mid W_i) + \ln P(W_i)
    \label{eq:gi_log}
\end{equation}

Expandindo a Eq.~\eqref{eq:gi_log} com a gaussiana $d$-dimensional
$\mathcal{N}(\boldsymbol{\mu}_i, \Sigma_i)$:

\begin{equation}
    g_i(\mathbf{x}) =
        -\frac{1}{2}(\mathbf{x}-\boldsymbol{\mu}_i)^T
         \Sigma_i^{-1}(\mathbf{x}-\boldsymbol{\mu}_i)
        + \ln P(W_i) + c_i
    \label{eq:gi_quadratico}
\end{equation}

\noindent onde $c_i = -\tfrac{d}{2}\ln 2\pi - \tfrac{1}{2}\ln|\Sigma_i|$
é a constante normalizadora. Este é o \textbf{discriminante
quadrático geral}, pois a superfície de decisão
$g_i(\mathbf{x}) - g_j(\mathbf{x}) = 0$ é uma quádrica (elipse,
hipérbole, par de retas).

Expandindo matricialmente:

\begin{equation}
    g_i(\mathbf{x}) =
        -\frac{1}{2}\mathbf{x}^T\Sigma_i^{-1}\mathbf{x}
        +\frac{1}{2}\mathbf{x}^T\Sigma_i^{-1}\boldsymbol{\mu}_i
        -\frac{1}{2}\boldsymbol{\mu}_i^T\Sigma_i^{-1}\boldsymbol{\mu}_i
        +\frac{1}{2}\boldsymbol{\mu}_i^T\Sigma_i^{-1}\mathbf{x}
        + \ln P(W_i) + c_i
    \label{eq:gi_expandido}
\end{equation}

\subsection{Superfície de Decisão Linear --- $\Sigma_i = \Sigma$}

Quando todas as classes compartilham a \textbf{mesma} matriz de
covariância ($\Sigma_i = \Sigma$), o termo quadrático
$\mathbf{x}^T\Sigma_i^{-1}\mathbf{x}$ e a constante $c_i$ são
idênticos em todas as classes e se cancelam em
$g_i(\mathbf{x})-g_j(\mathbf{x})=0$. O discriminante se reduz a:

\begin{equation}
    g_i(\mathbf{x}) =
        \frac{1}{2}\mathbf{x}^T\Sigma^{-1}\boldsymbol{\mu}_i
        -\frac{1}{2}\boldsymbol{\mu}_i^T\Sigma^{-1}\boldsymbol{\mu}_i
        +\frac{1}{2}\boldsymbol{\mu}_i^T\Sigma^{-1}\mathbf{x}
        + \ln P(W_i)
    \label{eq:gi_linear}
\end{equation}

Definindo $\mathbf{w}_i = \Sigma^{-1}\boldsymbol{\mu}_i$ e
$w_{i0} = \ln P(W_i) - \tfrac{1}{2}\boldsymbol{\mu}_i^T\Sigma^{-1}\boldsymbol{\mu}_i$,
a Eq.~\eqref{eq:gi_linear} torna-se simplesmente:

\begin{equation}
    \boxed{g_i(\mathbf{x}) = \mathbf{w}_i^T\mathbf{x} + w_{i0}}
    \label{eq:lda}
\end{equation}

\noindent A superfície de decisão $g_i(\mathbf{x}) = g_j(\mathbf{x})$
é um \textbf{hiperplano}, ortogonal à transformação linear
$\Sigma^{-1}(\boldsymbol{\mu}_i - \boldsymbol{\mu}_j)$.

\subsection{Caso Especial --- $\Sigma = \sigma^2 I$}

Assumindo adicionalmente que todos os atributos são descorrelacionados
e possuem a mesma variância $\sigma^2$:

\begin{equation}
    g_i(\mathbf{x}) =
        \frac{1}{\sigma^2}\boldsymbol{\mu}_i^T\mathbf{x} + w_{i0}
    \label{eq:lda_esferico}
\end{equation}

O hiperplano de decisão passa pelo ponto médio
$\mathbf{x}_0 = \tfrac{1}{2}(\boldsymbol{\mu}_i + \boldsymbol{\mu}_j)$
(para classes equiprováveis) e é ortogonal ao vetor
$\boldsymbol{\mu}_i - \boldsymbol{\mu}_j$. Este caso equivale ao
classificador por distância mínima ao centroide generalizado.

\subsection{Quatro Situações da Matriz de Covariância}

A Tabela~\ref{tab:situacoes} resume os quatro modos implementados,
correspondendo à instrução do professor de
\textit{``verificar várias situações diferentes para as matrizes de
covariância das classes''}.

\begin{table}[H]
    \centering
    \caption{Modos de covariância implementados.}
    \label{tab:situacoes}
    \begin{tabular}{@{}llll@{}}
        \toprule
        \textbf{Modo} & \textbf{Covariância} & \textbf{Equação} & \textbf{Fronteira} \\
        \midrule
        QDA completo  & $\Sigma_i$ completa por classe  & Eq.~\eqref{eq:gi_quadratico} & Quádrica \\
        QDA diagonal  & $\Sigma_i$ diagonal por classe  & Eq.~\eqref{eq:gi_quadratico} & Quádrica (eixos alinhados) \\
        LDA \textit{pooled} & $\Sigma$ compartilhada (pooled) & Eq.~\eqref{eq:lda}          & Hiperplano \\
        LDA esférico  & $\Sigma = \sigma^2 I$            & Eq.~\eqref{eq:lda_esferico} & Hiperplano ortogonal \\
        \bottomrule
    \end{tabular}
\end{table}

\subsection{Classificadores de Referência}

\textbf{Bayesiano Multivariado} — mesmo discriminante da
Eq.~\eqref{eq:gi_quadratico}, implementado nos trabalhos anteriores.

\textbf{Naive Bayes} — verossimilhança como produto de gaussianas 1D,
ignorando correlações entre atributos.

\textbf{DMC} — classifica pelo centroide mais próximo:
$\hat{C} = \arg\min_{C_i}\|\mathbf{x}-\boldsymbol{c}_i\|_2$.

\textbf{KNN} ($k=5$) — voto majoritário entre os $k$ vizinhos mais
próximos no conjunto de treino.

\subsection{Escolha do Par de Atributos para Visualização}

O par de atributos foi selecionado pelo critério de Fisher:

\begin{equation}
    \text{score}(a_1, a_2) = \text{Fisher}(a_1) + \text{Fisher}(a_2),
    \quad
    \text{Fisher}(a) =
    \frac{\text{variância entre classes}(a)}{\text{variância dentro das classes}(a)}
\end{equation}

% =============================================================
\section{Conjuntos de Dados}
% =============================================================

\subsection{Flor Íris}
150 amostras, 4 atributos contínuos, 3 classes balanceadas (Setosa,
Versicolor e Virginica). Dataset limpo, sem valores faltantes.

\subsection{Coluna Vertebral}
310 amostras, 6 atributos biomecânicos contínuos, 3 classes
ortopédicas (Normal, Hérnia, Espondilolistese). Completamente
numérico.

\subsection{Breast Cancer Wisconsin}
569 amostras, 30 atributos numéricos, 2 classes (Benigno e Maligno).
Versão \textit{Diagnostic} (UCI id=17).

\subsection{Dermatology}
366 amostras, 34 atributos, 6 classes de doenças dermatológicas.
Valores faltantes imputados pela média da coluna.

\subsection{Artificial I}
Dataset sintético gerado com 3 classes gaussianas 2D
($n=50$ por classe, semente $42$):
\begin{itemize}
    \item \textbf{Círculo:} $\boldsymbol{\mu}=(1{,}0,\;4{,}0)$, $\sigma=0{,}5$
    \item \textbf{Triângulo:} $\boldsymbol{\mu}=(4{,}0,\;4{,}0)$, $\sigma=0{,}5$
    \item \textbf{Estrela:} $\boldsymbol{\mu}=(2{,}0,\;2{,}0)$, $\sigma=0{,}5$
\end{itemize}

% =============================================================
\section{Metodologia}
% =============================================================

Foram realizadas \textbf{25 realizações} independentes por dataset,
com divisão 80\% treino e 20\% teste (embaralhamento aleatório a cada
realização). Para cada dataset foram executados os \textbf{oito
classificadores}: Bayesiano Multivariado, Naive Bayes, DMC, KNN,
QDA~completo, QDA~diagonal, LDA~\textit{pooled} e LDA~esférico.

A matriz de confusão foi extraída da realização com acurácia
\textbf{mais próxima da média} do QDA completo, por ser o
classificador central deste trabalho. As superfícies de decisão foram
plotadas para o QDA~completo e o LDA~\textit{pooled}, sobre o par de
atributos selecionado pelo critério de Fisher, exibindo os dados de
treino~($\bullet$) e teste~($\times$).

A matriz de covariância \textit{pooled} foi estimada por:

\begin{equation}
    \hat{\Sigma} =
    \frac{\displaystyle\sum_{i=1}^{c}(n_i - 1)\,\hat{\Sigma}_i}{N - c}
\end{equation}

\noindent onde $n_i$ é o número de amostras da classe $i$, $N$ é o
total de amostras de treino e $c$ é o número de classes. Uma
regularização $\epsilon I$ ($\epsilon = 10^{-6}$) foi aplicada antes
de inverter qualquer matriz de covariância, para evitar singularidades.

% =============================================================
\section{Resultados e Discussão}
% =============================================================

\subsection{Acurácia Média e Desvio Padrão}

\begin{table}[H]
    \centering
    \caption{Resultados dos trabalhos anteriores --- 25 realizações, divisão 80/20.}
    \label{tab:results_anteriores}
    \begin{tabular}{@{}lcccc@{}}
        \toprule
        \textbf{Base de Dados} & \textbf{Bayes. Mult.} & \textbf{Naive Bayes} & \textbf{KNN ($k{=}5$)} & \textbf{DMC} \\
        \midrule
        Íris             & $97{,}5 \pm 3{,}0$ & $94{,}8 \pm 4{,}3$ & $96{,}3 \pm 2{,}9$ & $90{,}9 \pm 4{,}9$ \\
        Artificial I     & $99{,}2 \pm 1{,}4$ & $99{,}5 \pm 1{,}2$ & $99{,}2 \pm 1{,}7$ & $99{,}2 \pm 1{,}4$ \\
        Coluna Vertebral & $83{,}9 \pm 4{,}5$ & $83{,}0 \pm 4{,}4$ & $82{,}3 \pm 3{,}3$ & $75{,}7 \pm 5{,}5$ \\
        Breast Cancer    & $88{,}4 \pm 3{,}0$ & $93{,}8 \pm 1{,}9$ & $93{,}6 \pm 2{,}2$ & $87{,}9 \pm 2{,}7$ \\
        Dermatology      & $88{,}6 \pm 3{,}8$ & $84{,}9 \pm 4{,}4$ & $86{,}2 \pm 2{,}9$ & $50{,}4 \pm 6{,}2$ \\
        \bottomrule
    \end{tabular}
\end{table}

\begin{table}[H]
    \centering
    \caption{Resultados do Trabalho~04 (LDA/QDA) --- 25 realizações, divisão 80/20.}
    \label{tab:results_ldaqda}
    \begin{tabular}{@{}lcccc@{}}
        \toprule
        \textbf{Base de Dados}
            & \textbf{QDA completo}
            & \textbf{QDA diagonal}
            & \textbf{LDA \textit{pooled}}
            & \textbf{LDA esférico} \\
        \midrule
        Íris             & $\cdot \pm \cdot$ & $\cdot \pm \cdot$ & $\cdot \pm \cdot$ & $\cdot \pm \cdot$ \\
        Artificial I     & $\cdot \pm \cdot$ & $\cdot \pm \cdot$ & $\cdot \pm \cdot$ & $\cdot \pm \cdot$ \\
        Coluna Vertebral & $\cdot \pm \cdot$ & $\cdot \pm \cdot$ & $\cdot \pm \cdot$ & $\cdot \pm \cdot$ \\
        Breast Cancer    & $\cdot \pm \cdot$ & $\cdot \pm \cdot$ & $\cdot \pm \cdot$ & $\cdot \pm \cdot$ \\
        Dermatology      & $\cdot \pm \cdot$ & $\cdot \pm \cdot$ & $\cdot \pm \cdot$ & $\cdot \pm \cdot$ \\
        \bottomrule
    \end{tabular}
\end{table}

\noindent\textit{Nota: os valores de ``$\cdot$'' devem ser
substituídos pelos resultados obtidos na execução do código.}

\subsection{Matrizes de Confusão}

Matrizes de confusão geradas pelo QDA completo na realização com
taxa de acerto mais próxima da média, por ser o classificador central
deste trabalho.

\begin{figure}[H]
    \centering
    \includegraphics[width=0.7\textwidth]{iris-matrix-confusao-qda.png}
    \caption{Matriz de confusão --- Íris (QDA completo).}
    \label{fig:conf_iris}
\end{figure}

\begin{figure}[H]
    \centering
    \includegraphics[width=0.7\textwidth]{vertebral-matrix-confusao-qda.png}
    \caption{Matriz de confusão --- Vertebral Column (QDA completo).}
    \label{fig:conf_vertebral}
\end{figure}

\begin{figure}[H]
    \centering
    \includegraphics[width=0.5\textwidth]{breast-matrix-confusao-qda.png}
    \caption{Matriz de confusão --- Breast Cancer (QDA completo).}
    \label{fig:conf_breast}
\end{figure}

\begin{figure}[H]
    \centering
    \includegraphics[width=0.7\textwidth]{dermatology-matrix-confusao-qda.png}
    \caption{Matriz de confusão --- Dermatology (QDA completo).}
    \label{fig:conf_dermatology}
\end{figure}

\begin{figure}[H]
    \centering
    \includegraphics[width=0.7\textwidth]{artificial-matrix-confusao-qda.png}
    \caption{Matriz de confusão --- Artificial I (QDA completo).}
    \label{fig:conf_artificial}
\end{figure}

\subsection{Superfícies de Decisão}

Superfícies de decisão comparando QDA completo (fronteira quádrica)
e LDA \textit{pooled} (fronteira linear) para os datasets Artificial~I,
Íris e Coluna Vertebral, com dados de treino~($\bullet$) e
teste~($\times$).

\begin{figure}[H]
    \centering
    \includegraphics[width=1\textwidth]{artificial-superficie-qda.png}
    \caption{Superfície de decisão --- Artificial I (QDA completo).}
    \label{fig:sup_artificial_qda}
\end{figure}

\begin{figure}[H]
    \centering
    \includegraphics[width=1\textwidth]{artificial-superficie-lda.png}
    \caption{Superfície de decisão --- Artificial I (LDA \textit{pooled}).}
    \label{fig:sup_artificial_lda}
\end{figure}

\begin{figure}[H]
    \centering
    \includegraphics[width=1\textwidth]{iris-superficie-qda.png}
    \caption{Superfície de decisão --- Íris (QDA completo).}
    \label{fig:sup_iris_qda}
\end{figure}

\begin{figure}[H]
    \centering
    \includegraphics[width=1\textwidth]{iris-superficie-lda.png}
    \caption{Superfície de decisão --- Íris (LDA \textit{pooled}).}
    \label{fig:sup_iris_lda}
\end{figure}

\begin{figure}[H]
    \centering
    \includegraphics[width=1\textwidth]{vertebral-superficie-qda.png}
    \caption{Superfície de decisão --- Coluna Vertebral (QDA completo).}
    \label{fig:sup_vertebral_qda}
\end{figure}

\begin{figure}[H]
    \centering
    \includegraphics[width=1\textwidth]{vertebral-superficie-lda.png}
    \caption{Superfície de decisão --- Coluna Vertebral (LDA \textit{pooled}).}
    \label{fig:sup_vertebral_lda}
\end{figure}

\subsection{Discussão}

\subsubsection{QDA completo vs.\ QDA diagonal}

O QDA com matriz completa captura as correlações entre atributos de
cada classe, modelando fronteiras elípticas orientadas livremente no
espaço de características. O QDA diagonal assume que os atributos são
condicionalmente independentes por classe (covariâncias cruzadas
nulas), produzindo elipses com eixos alinhados aos atributos. Para
datasets onde as correlações intraclasse são relevantes — como a
Coluna Vertebral — espera-se que o QDA completo supere o diagonal.
Nos datasets de baixa dimensionalidade como o Artificial~I (gerado com
covariância isotrópica), ambos devem apresentar desempenho similar.

\subsubsection{LDA \textit{pooled} vs.\ LDA esférico}

O LDA \textit{pooled} estima uma matriz de covariância única
$\hat{\Sigma}$ a partir de todas as classes, preservando informações
de espalhamento e correlação globais. O LDA esférico restringe ainda
mais: $\Sigma = \sigma^2 I$, assumindo variâncias iguais e atributos
descorrelacionados globalmente. Quando essa suposição é válida, ambos
produzem fronteiras lineares de desempenho equivalente; quando não,
o LDA \textit{pooled} tem vantagem.

\subsubsection{LDA vs.\ QDA}

O LDA impõe a restrição de fronteira linear (hiperplano), o que reduz
a variância do modelo mas aumenta o viés quando as classes realmente
possuem covariâncias distintas. O QDA é mais flexível (fronteiras
quádricas) e tende a superá-lo em datasets com classes de formas muito
diferentes, porém requer a estimativa de $c$ matrizes $d \times d$
separadas — o que pode ser ruidoso com poucas amostras por classe.
No Breast~Cancer ($d=30$), por exemplo, espera-se que o LDA leve
vantagem por evitar overfitting na estimativa de $\Sigma_i$.

\subsubsection{Comparação com trabalhos anteriores}

O Bayesiano Multivariado dos trabalhos anteriores é essencialmente o
mesmo que o QDA completo aqui implementado (a diferença está apenas na
forma de computar a classificação: \textit{posteriori} explícita
versus discriminante log-linear). Portanto, os resultados devem
concordar numericamente para validar a implementação.

% =============================================================
\section{Conclusão}
% =============================================================

Os quatro modos de covariância implementados formam uma hierarquia de
flexibilidade versus custo computacional/amostral: QDA~completo $>$
QDA~diagonal $>$ LDA~\textit{pooled} $>$ LDA~esférico. O modelo mais
adequado depende do dataset: em alta dimensionalidade com poucas
amostras (Breast~Cancer, Dermatology), modelos mais restritivos
tendem a generalizar melhor; em baixa dimensionalidade com classes bem
separadas (Artificial~I), todos convertem para resultados similares.

A principal contribuição deste trabalho frente aos anteriores é a
demonstração de que fronteiras lineares (LDA) são suficientes quando
a hipótese de covariância compartilhada é razoável, representando um
ganho direto em interpretabilidade e em custo de estimativa de
parâmetros, sem perda significativa de acurácia em boa parte dos
datasets avaliados.

\newpage
\begin{thebibliography}{9}
    \bibitem{duda} DUDA, R. O.; HART, P. E.; STORK, D. G\@.
        \textbf{Pattern Classification}. 2.~ed. Wiley-Interscience, 2001.
    \bibitem{bishop} BISHOP, C. M\@.
        \textbf{Pattern Recognition and Machine Learning}. Springer, 2006.
    \bibitem{slides} ROCHA NETO, A\@.
        \textbf{Densidade de Probabilidade Gaussiana} --- Notas de Aula.
        PPGET/IFCE, 2013.
    \bibitem{uci} DUA, D.; GRAFF, C\@. UCI Machine Learning Repository.
        Disponível em: \url{https://archive.ics.uci.edu/ml}. Acesso em: \today.
\end{thebibliography}

\end{document}